\section{Closing Remarks}
\label{sec:closing}

The real world does not respect the boundaries of an agent's training data. Objects appear that were never labeled, dynamics change in ways that were never simulated, and tasks evolve beyond what was anticipated during system design. The question is not whether agents will encounter novelty, but whether they can respond to it effectively.

We have argued that the answer depends on how agents represent their interaction with the world. Monolithic representations that tightly couple perception, reasoning, and action provide no natural interfaces for diagnosing failures or localizing adaptation. Structured abstractions---representations that decompose behavior into compositional sequences and organize interaction around task-relevant invariants---create these interfaces. They enable agents to identify what has changed, restrict learning to what must change, and preserve what still works.

The formal frameworks, benchmark environments, adaptation methods, and generalization techniques developed in this dissertation provide concrete instantiations of this principle. They show that structured abstractions over actions make accommodation after execution failure more localized and efficient, and that structured abstractions over interaction expand the variation that learned skills can assimilate through generalization across objects, physical parameters, and robotic platforms. Beyond the specific algorithms and systems, the formal vocabulary itself---novelty as agent-relative mismatch, the taxonomies by task impact and domain element, the integrated planning task, and the distinction between action and interaction abstractions---constitutes a contribution that we hope will be useful to others working on open-world autonomy.

Together, these contributions offer a perspective on open-world autonomy in which the structure of an agent's representations is not merely a design choice but a fundamental determinant of its capacity to operate in a world that is always, in some respect, new.
